% GENERATED FILE. Edit canonical lessons, then run npm run generate:books.
% canonical-source-hash: fnv1a64:a960739973bbbc70
% canonical-lessons: ML-C382-kalapam, ML-C382-ayudham, ML-C382-valayuka-surround, ML-C382-turattuka, ML-C382-santamakuka

\chapter{A Riot in the Town}
\label{ch:ml-a-riot-in-the-town}
\hlchaptermodality{382}

\begin{chapteropening}
\textbf{By the end of this chapter:} \emph{I can say a riot, a weapon, to surround, to drive away and to calm down.}

Complete the last lesson of chapter 382: I can say a riot, a weapon, to surround, to drive away and to calm down.
\end{chapteropening}

\section[\texorpdfstring{\ml{കലാപം} (kalāpaṁ)}{കലാപം (kalāpaṁ)}]{\ml{കലാപം} (kalāpaṁ) — a riot}

\label{lesson:ML-C382-kalapam}



\begin{quote}
Before the new one: say the Malayalam for to whisper, then the Malayalam for to disappear.
\end{quote}

\subsection*{You'll want to know: \ml{കലാപം}}
\textbf{\ml{കലാപം}} — \emph{kalāpaṁ} — \textquotedblleft{}a riot\textquotedblright{}.

\begin{quote}
\ml{കലാപം} \ml{കാരണം} \ml{കടകൾ} \ml{അടച്ചു}.
\end{quote}

\begin{quote}
\emph{kalāpaṁ kāraṇaṁ kaṭakaḷ aṭaccu.}
\end{quote}

\begin{quote}
Shops closed because of the riot.
\end{quote}

Also an uproar at a meeting.

\begin{cousinweb}
The origin of \textbf{\ml{കലാപം}} (\emph{kalāpaṁ}) in the sense of a riot is not settled. In the news it is violent disorder in a town.
\end{cousinweb}

\begin{grammarlens}[title={What you've built}]
One more word for telling what happened.
\end{grammarlens}

\subsection*{Your turn}
\begin{itemize}
\raggedright
  \item \emph{Say it:} \emph{kalāpaṁ}
  \item \emph{Say it:} \emph{kalāpaṁ}, once more
  \item \emph{Say it:} \emph{kalāpaṁ kāraṇaṁ kaṭakaḷ aṭaccu.} — Shops closed because of the riot.
  \item \emph{Recall:} say the Malayalam for to whisper, then the Malayalam for to disappear, then say \emph{kalāpaṁ} again
\end{itemize}

\begin{tcolorbox}[breakable,colback=teal!4,colframe=teal!35!black,title={Before you move on}]
What does \ml{കലാപം} mean? (\textquotedblleft{}A riot\textquotedblright{}.) Say it once more.
\end{tcolorbox}
\section[\texorpdfstring{\ml{ആയുധം} (āyudhaṁ)}{ആയുധം (āyudhaṁ)}]{\ml{ആയുധം} (āyudhaṁ) — a weapon}

\label{lesson:ML-C382-ayudham}



\begin{quote}
Before the new one: say the Malayalam for to disappear, then the Malayalam for a riot.
\end{quote}

\subsection*{You'll want to know: \ml{ആയുധം}}
\textbf{\ml{ആയുധം}} — \emph{āyudhaṁ} — \textquotedblleft{}a weapon\textquotedblright{}.

\begin{quote}
\ml{പോലീസ്} \ml{കള്ളന്റെ} \ml{കയ്യിൽ} \ml{നിന്ന്} \ml{ആയുധം} \ml{പിടിച്ചെടുത്തു}.
\end{quote}

\begin{quote}
\emph{pōlīsŭ kaḷḷanṟe kayyil ninnŭ āyudhaṁ piṭicceṭuttu.}
\end{quote}

\begin{quote}
The police seized a weapon from the robber's hand.
\end{quote}

Plural: \ml{ആയുധങ്ങൾ}, \emph{āyudhaṅṅaḷ}.

A tool is \ml{ഉപകരണം}; a weapon is \ml{ആയുധം}.

\begin{cousinweb}
\textbf{\ml{ആയുധം}} (\emph{āyudhaṁ}) is borrowed from Sanskrit \emph{āyudha}, \textquotedblleft{}a weapon\textquotedblright{} (CDIAL 1291), from the root \emph{yudh}, \textquotedblleft{}to fight\textquotedblright{}, the root of \ml{യുദ്ധം} (\emph{yuddhaṁ}).
\end{cousinweb}

\begin{grammarlens}[title={What you've built}]
One more word for telling what happened.
\end{grammarlens}

\subsection*{Your turn}
\begin{itemize}
\raggedright
  \item \emph{Say it:} \emph{āyudhaṁ}
  \item \emph{Say it:} \emph{āyudhaṁ}, once more
  \item \emph{Say it:} \emph{pōlīsŭ kaḷḷanṟe kayyil ninnŭ āyudhaṁ piṭicceṭuttu.} — The police seized a weapon from the robber's hand.
  \item \emph{Recall:} say the Malayalam for to disappear, then the Malayalam for a riot, then say \emph{āyudhaṁ} again
\end{itemize}

\begin{tcolorbox}[breakable,colback=teal!4,colframe=teal!35!black,title={Before you move on}]
What does \ml{ആയുധം} mean? (\textquotedblleft{}A weapon\textquotedblright{}.) Say it once more.
\end{tcolorbox}
\section[\texorpdfstring{\ml{വളയുക} (vaḷayuka)}{വളയുക (vaḷayuka)}]{\ml{വളയുക} (vaḷayuka) — to surround}

\label{lesson:ML-C382-valayuka-surround}



\begin{quote}
Before the new one: say the Malayalam for a riot, then the Malayalam for a weapon.
\end{quote}

\subsection*{You'll want to know: \ml{വളയുക}}
\textbf{\ml{വളയുക}} — \emph{vaḷayuka} — \textquotedblleft{}to surround\textquotedblright{}.

\begin{quote}
\ml{പോലീസ്} \ml{ആ} \ml{വീട്} \ml{വളഞ്ഞു}.
\end{quote}

\begin{quote}
\emph{pōlīsŭ ā vīṭŭ vaḷaññu.}
\end{quote}

\begin{quote}
The police surrounded the house.
\end{quote}

Past: \ml{വളഞ്ഞു}, \emph{vaḷaññu}.

The same form also means to bend, as in \ml{വളവ്}, a bend.

\begin{cousinweb}
\textbf{\ml{വളയുക}} (\emph{vaḷayuka}) is native Dravidian: DEDR 5313 lists Malayalam \emph{vaḷayuka}, \textquotedblleft{}to surround\textquotedblright{}, with \ml{വള}, \textquotedblleft{}a ring, a bangle\textquotedblright{}: to make a ring round something.
\end{cousinweb}

\begin{grammarlens}[title={What you've built}]
One more word for telling what happened.
\end{grammarlens}

\subsection*{Your turn}
\begin{itemize}
\raggedright
  \item \emph{Say it:} \emph{vaḷayuka}
  \item \emph{Say it:} \emph{vaḷayuka}, once more
  \item \emph{Say it:} \emph{pōlīsŭ ā vīṭŭ vaḷaññu.} — The police surrounded the house.
  \item \emph{Recall:} say the Malayalam for a riot, then the Malayalam for a weapon, then say \emph{vaḷayuka} again
\end{itemize}

\begin{tcolorbox}[breakable,colback=teal!4,colframe=teal!35!black,title={Before you move on}]
What does \ml{വളയുക} mean? (\textquotedblleft{}To surround\textquotedblright{}.) Say it once more.
\end{tcolorbox}
\section[\texorpdfstring{\ml{തുരത്തുക} (turattuka)}{തുരത്തുക (turattuka)}]{\ml{തുരത്തുക} (turattuka) — to drive away; to chase off}

\label{lesson:ML-C382-turattuka}



\begin{quote}
Before the new one: say the Malayalam for a weapon, then the Malayalam for to surround.
\end{quote}

\subsection*{You'll want to know: \ml{തുരത്തുക}}
\textbf{\ml{തുരത്തുക}} — \emph{turattuka} — \textquotedblleft{}to drive away; to chase off\textquotedblright{}.

\begin{quote}
\ml{നാട്ടുകാർ} \ml{ആനയെ} \ml{കാട്ടിലേക്ക്} \ml{തുരത്തി}.
\end{quote}

\begin{quote}
\emph{nāṭṭukār ānaye kāṭṭilēkkŭ turatti.}
\end{quote}

\begin{quote}
The locals drove the elephant back into the forest.
\end{quote}

Past: \ml{തുരത്തി}, \emph{turatti}.

Of animals and of enemies.

\begin{cousinweb}
\textbf{\ml{തുരത്തുക}} (\emph{turattuka}) is native Dravidian: DEDR 3340 lists Malayalam \emph{turattuka}, \textquotedblleft{}to drive away\textquotedblright{}, beside Tamil \emph{tura}, \textquotedblleft{}to drive, to beat away\textquotedblright{}.
\end{cousinweb}

\begin{grammarlens}[title={What you've built}]
One more word for telling what happened.
\end{grammarlens}

\subsection*{Your turn}
\begin{itemize}
\raggedright
  \item \emph{Say it:} \emph{turattuka}
  \item \emph{Say it:} \emph{turattuka}, once more
  \item \emph{Say it:} \emph{nāṭṭukār ānaye kāṭṭilēkkŭ turatti.} — The locals drove the elephant back into the forest.
  \item \emph{Recall:} say the Malayalam for a weapon, then the Malayalam for to surround, then say \emph{turattuka} again
\end{itemize}

\begin{tcolorbox}[breakable,colback=teal!4,colframe=teal!35!black,title={Before you move on}]
What does \ml{തുരത്തുക} mean? (\textquotedblleft{}To drive away; to chase off\textquotedblright{}.) Say it once more.
\end{tcolorbox}
\section[\texorpdfstring{\ml{ശാന്തമാകുക} (śāntamākuka)}{ശാന്തമാകുക (śāntamākuka)}]{\ml{ശാന്തമാകുക} (śāntamākuka) — to calm down}

\label{lesson:ML-C382-santamakuka}



\begin{quote}
Before the new one: say the Malayalam for to surround, then the Malayalam for to drive away.
\end{quote}

\subsection*{You'll want to know: \ml{ശാന്തമാകുക}}
\textbf{\ml{ശാന്തമാകുക}} — \emph{śāntamākuka} — \textquotedblleft{}to calm down\textquotedblright{}.

\begin{quote}
\ml{കുറച്ച്} \ml{കഴിഞ്ഞപ്പോൾ} \ml{നഗരം} \ml{ശാന്തമായി}.
\end{quote}

\begin{quote}
\emph{kuṟaccŭ kaḻiññappōḷ nagaraṁ śāntamāyi.}
\end{quote}

\begin{quote}
After a while the city calmed down.
\end{quote}

Past: \ml{ശാന്തമായി}, \emph{śāntamāyi}.

\ml{ശാന്തം}, calm, + \ml{ആകുക}, to become.

\begin{cousinweb}
\textbf{\ml{ശാന്തമാകുക}} (\emph{śāntamākuka}) is built from Sanskrit \emph{śānta}, \textquotedblleft{}abated, calm\textquotedblright{} (CDIAL 12390), borrowed as \ml{ശാന്തം}, and the Malayalam \ml{ആകുക}, \textquotedblleft{}to become\textquotedblright{}.
\end{cousinweb}

\begin{grammarlens}[title={What you've built}]
One more word for telling what happened.
\end{grammarlens}

\subsection*{Your turn}
\begin{itemize}
\raggedright
  \item \emph{Say it:} \emph{śāntamākuka}
  \item \emph{Say it:} \emph{śāntamākuka}, once more
  \item \emph{Say it:} \emph{kuṟaccŭ kaḻiññappōḷ nagaraṁ śāntamāyi.} — After a while the city calmed down.
  \item \emph{Recall:} say the Malayalam for to surround, then the Malayalam for to drive away, then say \emph{śāntamākuka} again
\end{itemize}

\begin{tcolorbox}[breakable,colback=teal!4,colframe=teal!35!black,title={Before you move on}]
What does \ml{ശാന്തമാകുക} mean? (\textquotedblleft{}To calm down\textquotedblright{}.) Say it once more.
\end{tcolorbox}
